\chapter{Motif Statistics and Analytic Information Theory}
\label{ch17}

This chapter points toward a mature classical toolkit rather than trying to redevelop it from first principles. Analytic information theory studies sequence statistics --- motif counts, trie depth, match lengths, compression parameters --- by encoding them into finite-state or transfer-matrix objects and then applying complex-analytic tools such as singularity analysis, Mellin transforms, and saddle-point methods.

The exact theorems belong to classical finite-state sources. The possible transfer to WFA approximations of LLMs is a research program, not an automatic corollary. The chapter will keep that distinction explicit.

\section{A Classical Theme}

For finite-state or Markov sources, many statistics of interest can be packaged into generating functions or matrix resolvents. Once that is done, asymptotic analysis can often be read off from poles, perturbations of dominant eigenvalues, Mellin poles, or saddle points. That is why analytic information theory fits so naturally into the rest of this book: it is the information-theoretic analogue of analytic combinatorics.

\section{Motif Statistics: What Is the Generating Function?}

Fix a source on strings over an alphabet $\Sigma$, and fix a motif $M$ (a word or a regular-language pattern). There are two related but different generating functions one can study.

\begin{itemize}
  \item A \textbf{counting} generating function that counts how many strings of length $n$ have exactly $k$ occurrences of the motif.
  \item A \textbf{weighted / probabilistic} generating function in which the coefficient is the probability that a random length-$n$ output has exactly $k$ occurrences.
\end{itemize}

For analytic information theory, the second is often the more natural object. We therefore define
\[
  F(z,u) = \sum_{n \ge 0}\sum_{k \ge 0} p_{n,k} z^n u^k,
\]
where
\[
  p_{n,k} = \Pr(\text{a random length-}n\text{ string has exactly }k\text{ occurrences of }M).
\]
If one wants the unweighted counting version instead, one simply replaces probabilities by counts. The same automaton technology underlies both.

\section{Autocorrelation and Overlap}

The structure of motif occurrences is governed by overlap.

\begin{definition}[Autocorrelation polynomial]
If
\[
  M = m_1m_2\cdots m_\ell,
\]
then its autocorrelation polynomial is
\[
  C_M(z) = \sum_{k=0}^{\ell-1} c_k z^k,
\]
where $c_k=1$ if the suffix of $M$ of length $k$ equals the prefix of $M$ of length $k$, and $c_k=0$ otherwise.
\end{definition}

This polynomial records how one occurrence of the motif can overlap with the next. It matters because overlaps create dependence between nearby occurrence indicators.

\begin{example}[The motif ``aba'']
For
\[
  M=\texttt{aba},
\]
there is always the trivial overlap $k=0$, and there is also a 1-letter overlap because the suffix `a' equals the prefix `a'. But there is no 2-letter overlap because `ba' differs from `ab'. Thus
\[
  C_M(z)=1+z.
\]

a pattern with overlap can occur in strings like `ababa', where two occurrences share one letter.
\end{example}

\section{The Automaton Recipe}

The usual workflow is mechanical.

\begin{enumerate}
  \item Build a finite automaton whose states record how much of the motif has currently been matched (an overlap automaton, essentially the finite-state suffix-matching machine behind pattern search).
  \item Add a mark $u$ whenever a transition completes a new occurrence.
  \item Weight each transition by $z$ for length, and by the source probability if a probabilistic source is being used.
  \item Form the transfer matrix $T(z,u)$ of this weighted automaton.
\end{enumerate}

Then the bivariate generating function has the rational form
\[
  F(z,u) = \mathbf{v}^T (I-T(z,u))^{-1} \mathbf{e}
\]
for appropriate initial and terminal vectors.

That is the finite-state miracle: once the pattern-tracking automaton is built, the rest is linear algebra.

\begin{theorem}[Gaussian limit law for motif counts, black box]
For a fixed motif tracked by a finite automaton under an ergodic Markov source, the number of motif occurrences in a random length-$n$ string has mean
\[
  \mu n + O(1),
\]
variance
\[
  \sigma^2 n + O(1),
\]
and a central-limit theorem after centering and scaling.
\end{theorem}

The mechanism is singularity perturbation in the auxiliary variable $u$ together with a quasi-power theorem. For this chapter, the safe moral is enough: finite-state motif tracking leads naturally to linear growth of mean and variance and often to a Gaussian limit law.

\section{A Simple ``aba'' Mean Calculation}

For the uniform i.i.d. binary source and the motif `aba', the leading mean can be obtained without any heavy machinery. Let
\[
  X_n = \#\text{occurrences of ``aba'' in a random binary string of length }n.
\]
There are $n-2$ possible starting positions, and at each one the probability of seeing `aba' is
\[
  \left(\frac12\right)^3 = \frac18.
\]
So by linearity of expectation,
\[
  \mathbb E[X_n] = \frac{n-2}{8}.
\]
Overlap matters for the variance and the full distribution, but not for this first-moment calculation.

\section{Trie Depth and Entropy}

Tries give another classical example. Build a trie from independent strings drawn from a finite-state source. The depth needed to distinguish one string from the others is closely tied to the source entropy rate.

\begin{theorem}[Trie depth under a Markov source, black box]
If strings are generated by an ergodic Markov source with Shannon entropy rate $h$, then the expected depth of a random key in a trie built from $n$ independent samples satisfies
\[
  \mathbb E[D_n] \sim \frac{1}{h}\log n.
\]
A finer expansion often includes a bounded periodic fluctuation in $\log n$.
\end{theorem}

The important point is conceptual: the constant is the reciprocal of the \emph{Shannon entropy rate of the source itself}, not a naive path-counting growth exponent. Chapter~9 already warned that support growth and Shannon entropy are different objects, and that warning remains in force here.

For the uniform binary source,
\[
  h=\log 2,
\]
so
\[
  \mathbb E[D_n] \sim \log_2 n.
\]
That is exactly what one should expect: among $n$ roughly equally likely binary strings, one needs about $\log_2 n$ bits to isolate a typical one.

\section{Two Classical Analytic Tools}

Two techniques appear repeatedly in this literature.

\paragraph{Mellin transforms.}
Mellin methods are especially good for digital trees and periodic fluctuations. Non-real poles in Mellin space translate into oscillatory terms in $\log n$.

\paragraph{Saddle-point methods.}
When the relevant generating function is entire rather than singular at a finite point, one often evaluates coefficients by locating the saddle point of the Cauchy integral. This is an advanced tool and we use it here only as signposting.

A student should not feel compelled to master both techniques from this chapter alone. The important lesson is structural: different statistics call for different analytic machines, but the machines are reusable once the right finite-state encoding has been built.

\section{What Transfers to WFA Approximations?}

This is the chapter's main cautionary question.

If a WFA surrogate for an LLM preserves the statistic of interest to sufficient accuracy, then the exact finite-state theorems above become plausible approximations for the surrogate. But that ``if'' hides the hard part.

For example:
\begin{itemize}
  \item a surrogate that matches next-token conditionals on short contexts may still get long-range motif counts wrong;
  \item a surrogate that matches the length PGF may still get phrase statistics wrong;
  \item a surrogate with excellent finite-length agreement may still have the wrong dominant poles asymptotically.
\end{itemize}

So the safe formulation is this:
\begin{quote}
Analytic information theory supplies an exact toolkit for exact finite-state models and a research program for sufficiently faithful finite-state approximations.
\end{quote}
It does not supply a theorem that every WFA approximation of an LLM automatically preserves every asymptotic statistic we care about.

\section*{Summary}

This chapter's clean message is:
\begin{enumerate}
  \item finite-state encodings of sequence statistics lead to rational or transfer-matrix generating functions;
  \item once those encodings are in hand, classical analytic tools yield sharp asymptotics;
  \item transferring those results to LLM approximations is promising, but conditional on the quality and target of the approximation.
\end{enumerate}

That is enough to motivate the analytic information theory program without overselling what has already been proved for LLMs.

\section*{Further Reading}

\begin{itemize}
  \item \textbf{W. Szpankowski}, \emph{Average Case Analysis of Algorithms on Sequences} (2001) \cite{Szpankowski2001} --- the central monograph on analytic information theory.
  \item \textbf{L. J. Guibas and A. M. Odlyzko}, ``String overlaps, pattern matching, and nontransitive games'' (1981) \cite{GuibasOdlyzko1981} --- for autocorrelation and overlap structure.
  \item \textbf{P. Nicodème, B. Salvy, and P. Flajolet}, ``Motif statistics'' (2002) \cite{NicodemeSalvyFlajolet2002} --- for transfer-matrix motif counting and Gaussian laws.
  \item \textbf{P. Jacquet and W. Szpankowski}, ``Analysis of digital tries with Markovian dependency'' (1991) \cite{JacquetSzpankowski1991} --- for trie depth and Mellin techniques.
  \item \textbf{P. Flajolet and R. Sedgewick}, \emph{Analytic Combinatorics} (2009) \cite{FlajoletSedgewick2009} --- for quasi-powers, singularity analysis, and saddle-point methods.
\end{itemize}

\section*{Exercises}

\begin{exercise}
Compute the autocorrelation polynomial of the motif `abab'. Which nontrivial overlaps occur?
\end{exercise}

\begin{exercise}
Explain why the strict distinction between counting and weighted motif generating functions matters by giving one statistic that changes when source probabilities are introduced and one that does not.
\end{exercise}
